\section{Weighted interval orbits with polynomial output}
\label{sec:weighted-orbits}

Weighted orbit counting is classical. Agol, Hass, and Thurston prove a
polynomial-time weighted extension of their interval algorithm in Theorem~16;
their Corollary~17 applies it to the normal coordinates of connected surface
components \cite{AHT2006}. We use this antecedent explicitly. The contribution
here is a maintained implementation with a precise suffix-transport invariant,
a compressed histogram contract, independently checked arithmetic, and the
small set of weights needed by the subsequent disc query. Neither weighted
counting itself nor the deduction that normal components can be recovered in
polynomial time is claimed as new.

\subsection{Input, output, and the surviving-representative invariant}

Let $X=\{0,\ldots,N-1\}$, with an equivalence relation generated by $K$
interval pairings. Each pairing is an isometry between equal-length integer
intervals, with either sign. Endpoints and $N$ are binary integers. A weight
function $w:X\to\Z^d$ is supplied by $s$ additive half-open intervals:
\[
 w(x)=\sum_{i=1}^{s} v_i\,\mathbf 1_{[\ell_i,r_i)}(x),
 \qquad v_i\in\Z^d.
\]
Intervals may overlap, coordinates may be negative, and uncovered points have
zero weight. Sorting their endpoints and summing differences produces a
canonical list of constant runs covering $X$. Adjacent equal runs are merged.
For $N>0$ the initial number $R_0$ of runs is at most $2s+1$; the empty universe
has no runs.

For an orbit $O$, define $W(O)=\sum_{x\in O}w(x)$. The output is the sorted
histogram
\[
 H(v)=\#\{O:W(O)=v\}, \qquad H(v)>0
\]
on its explicitly represented support. Multiplicities remain binary integers.
The algorithm does not emit one record per orbit. For example, an unpaired
universe of size $2^b$ with constant weight $v$ has the one-record answer
$H(v)=2^b$. Expanding that answer into a list would destroy the claimed encoded
complexity. Distinct components with the same weight deliberately share a
record; their identities and attaching maps are not part of this contract.

The unweighted producer supplies a complete AHT reduction trace. It contains
identity deletion, reflection trimming, periodic merging, transmission,
static-interval contraction, and suffix truncation. Both the classical and
Fine--Wilf merger versions are supported by the existing independent trace
checker. Weight propagation reconstructs the current pairing list from these
events, without running the discovery scheduler again.

The relevant invariant is stronger than conservation of the total weight.
Each surviving point represents a disjoint collection of original points;
its current weight is their vector sum. The current orbit relation groups
these collections into exactly the original orbits not yet emitted. Identity
deletion, trimming, merging, and transmission change the generators of the
relation without removing points. They therefore leave all point weights
unchanged. Only the two point-removal operations require weight updates.

\subsection{Exact transport through a truncation}

Suppose the current universe is $[0,N)$ and a checked event removes the suffix
$[c,N)$. That suffix is incident only to its designated carrier pairing.
For a translation carrier with displacement $p>0$, its domain begins at $a$,
its range begins at $a+p$, and $c\ge a+p$. Repeated inverse application sends
every deleted point to the retained interval $J=[c-p,c)$ by
\begin{equation}
 f(x)=c-p+\bigl((x-(c-p))\bmod p\bigr).
 \label{eq:weighted-residue-map}
\end{equation}
All intermediate inverse applications are defined: until the first value
below $c$, the point remains in the carrier's range. The target belongs to
$[a,c)$ and is retained. Each original orbit still has a representative.

Consider one constant suffix run $[\ell,r)$ of value $v$. Write
\[
 r-\ell=qp+u,\qquad 0\le u<p,
 \qquad z=c-p+((\ell-(c-p))\bmod p).
\]
Its pushforward contributes $qv$ to every point of $J$. The remaining $u$
points contribute $v$ on the cyclic interval of length $u$ starting at $z$.
Putting $u_1=\min(u,c-z)$, this is the sum of contributions on
\[
 [z,z+u_1),\qquad [c-p,c-p+u-u_1).
\]
Empty intervals are omitted. Thus a run of astronomical length needs one
division with remainder and at most three constant-interval additions. The
retained prefix keeps its existing weights, and all suffix contributions
are added to it using an endpoint-difference table.

A checked reflection carrier has disjoint domain and range after trimming.
With current right endpoint $N-1$ and left endpoint $a$, its inverse is
$x\mapsto a+N-1-x$. Consequently a deleted run $[\ell,r)$ maps bijectively to
\[
 [a+N-r,a+N-\ell).
\]
This interval lies wholly in the retained prefix. The implementation adds
the same vector there and discards the source suffix. The reflection's
order reversal must be retained; treating every pairing as a translation
would give incorrect component sums even when the unweighted count agreed.

For a static contraction, every point in each supplied gap is a singleton
orbit of the current relation. Its accumulated weight is already the full
weight of one original orbit. Intersecting the gaps with the current runs
therefore adds the intersection lengths to the corresponding histogram
entries. The remaining points and weights are shifted by the same
order-preserving gap-removal map used for the pairings.

\begin{proposition}[Correct weighted replay]
\label{prop:weighted-correct}
Replaying these operations along a valid complete AHT trace returns exactly
the histogram of original orbit-weight sums, including signed weights.
\end{proposition}
\begin{proof}
Maintain the disjoint representative collections described above. The four
relation-only operations preserve them unchanged. A suffix truncation assigns
each deleted collection to its image under the carrier's inverse, possibly
iterated. These assignments partition all deleted collections and preserve
their original orbit membership. Summing their vectors gives exactly the
displayed pushforward formulas. A static point then represents one complete
original orbit and can be emitted. The checked trace finishes with no points
and no pairings, so every original orbit is emitted exactly once. No step uses
positivity; the proof applies coordinatewise in the additive group $\Z^d$.
\end{proof}

\subsection{A bound on run growth and histogram size}

The fact that one source run gives constantly many intervals is insufficient
to prove polynomial space: repeated constant-factor growth could still be
exponential. The necessary bound follows from transporting jumps.

\begin{lemma}[Additive run growth]
\label{lem:weighted-run-growth}
A truncation increases the number of canonical weight runs by at most two.
A static contraction does not increase it. Thus after $E$ trace events,
\[
 R\le R_*:=R_0+2E.
\]
\end{lemma}
\begin{proof}
Split the old runs at the truncation point into $P$ nonempty prefix runs and
$Q$ nonempty suffix runs. At most one old run was split, so
$P+Q\le R+1$. Extend the suffix function by zero to the integer line. Its
discrete derivative has jumps only at its $Q+1$ endpoints. Under translation
folding, each jump has one residue image modulo $p$. The cutoff endpoint
maps to the wrap boundary $c-p$, leaving at most $Q$ other interior residue
locations. Relative to the $P-1$ internal prefix boundaries, these locations
and the possible new wrap boundary give at most $P+Q+1$ runs. This is at most
$R+2$. Coincident endpoints and cancelling jumps only reduce the count.

For a reflection, the $Q+1$ suffix endpoints have one reflected image each.
Together with the $P-1$ prefix boundaries they again give at most
$P+Q+1\le R+2$ runs. Finally, deleting static intervals retains a subsequence
of the old constant runs. Every new change of value corresponds to at least
one old change between consecutive retained points. Coalescing adjacent equal
values therefore cannot increase the run count.
\end{proof}

If a contraction contains $g$ gaps, their intersections with a run list have
at most $R+g-1$ nonempty pieces. There are at most $2K+1$ gaps, since all
pairing domains and ranges together have at most $2K$ occupied intervals.
Pairing count never increases during the trace. Consequently the number $h$
of distinct histogram records satisfies the conservative bound
\begin{equation}
 h\le E\bigl(R_0+2E+2K+1\bigr).
 \label{eq:weighted-output-bound}
\end{equation}
Repeated equal weights share a record, so this bounds the actual output as
well as the number of possible emissions. Since the maintained AHT trace
length is polynomial in $K$ and $\log(N+1)$, the histogram has polynomial
encoded size. This observation is essential when the number of components
itself is exponentially large.

\subsection{Arithmetic and certificate complexity}

Assume every coordinate of every supplied $v_i$ has absolute value at most
$A$. Each current representative contains a subset of original points, hence
every accumulated coordinate has magnitude at most $NsA$. This also bounds
each completed orbit sum. Taking
\[
 B=O\bigl(1+\log(N+1)+\log(s+1)+\log(A+1)\bigr)
\]
therefore bounds all weight bit lengths. Multiplicities use
$O(1+\log(N+1))$ bits. Quotient multiplication is integer multiplication on
these bit lengths, not a unit-cost operation on arbitrarily large numbers.

Initial normalization uses $O(sd)$ coordinate additions and $O(s\log(s+1))$
endpoint comparisons. Given a trace of $E$ events, one truncation uses
$O(dR_*)$ coordinate operations and $O(R_*\log(R_*+1))$ endpoint comparisons.
A contraction requires at most
$O(d(R_*+K)+(R_*+K)\log(R_*+K+1))$ such operations. Source reconstruction and
the final histogram sort add polynomial costs; lexicographic sorting of $h$
vectors uses at most $O(dh\log(h+1))$ coordinate comparisons. A useful
event-sensitive accounting bound, apart from the existing unweighted
producer and checker, is therefore
\[
 O\!\left(sd+s\log(s+1)
 +E\bigl[d(R_*+K)+(R_*+K)\log(R_*+K+1)+K^2\bigr]
 +dh\log(h+1)\right).
\]
The $K^2$ term charges the independent weighted checker's direct reconstruction of contracted pairing endpoints across up to $O(K)$ gaps. The producer can use the sharper bound without that term. This counts arithmetic and ordinary dictionary operations. Integer hashing,
comparisons, additions, and products must then be charged their bit costs.
Adversarial collisions in Python dictionaries can worsen the displayed
container-operation model, but still give a polynomial bound for these
explicit polynomial-size tables; deterministic balanced maps give the usual
additional logarithmic alternative. No conclusion depends on treating a
$2^b$-bit integer as one machine word.

The certificate stores the canonical source weights, the existing unweighted
orbit proof, and the claimed sorted histogram. A contraction event can list
$O(K)$ gaps, so trace payload size is charged as well as event count. The run
bound, output bound, and bit bound together show polynomial certificate
length. Internally the producer creates an orbit trace even when the caller
does not request a returned certificate: weight propagation needs that
sequence of contractions. A supplied trace can instead be source-checked and
reused, avoiding a new orbit search.

\subsection{Independent arithmetic and sufficient statistics}

The independent weighted checker first verifies all unweighted local rules
against the caller's pairing list. It then reconstructs the pairing updates
and weight transport separately. In particular, its translation arithmetic
does not reuse the producer's quotient-and-remainder interval additions.
For the zero-extended suffix function $v$ and a target representative $y$, put
\[
 S(y)=\sum_{k\in\Z}v(y+kp).
\]
At the leftmost residue $y=c-p$, a source run $[\ell,r)$ contributes its
weight multiplied by
\[
 \left\lfloor\frac{r-1-y}{p}\right\rfloor
 -\left\lfloor\frac{\ell-1-y}{p}\right\rfloor.
\]
This determines one anchor value. The identity
\[
 S(y)-S(y-1)=\sum_{k\in\Z}\bigl(v(y+kp)-v(y-1+kp)\bigr)
\]
determines every remaining value by integrating the independently mapped
suffix jumps. These two arithmetic constructions reach the same required
pushforward through different calculations. The checker additionally checks
the total orbit count, exact sorted histogram, source binding, and signed
mass conservation. A partial histogram is never returned when the
underlying orbit allowance is exhausted; cancellation callbacks propagate
through construction and replay.

Finally, many decisions depend on a small additive signature rather than a
full component description. If an orbit predicate factors through $d_0$
additive statistics, those statistics may be transported directly and the
desired count obtained by summing histogram multiplicities satisfying the
predicate. The number of vector coordinates in the arithmetic bound becomes
$d_0$. In the next section, three statistics suffice for the disc predicate,
whereas full normal-component extraction uses $7t$ coordinates. This removes
a factor of order $t$ from the dimension-dependent arithmetic bound; it does
not remove common orbit-search and endpoint-processing costs. The signatures
also need not recover the normal vector of a particular disc. A subsequent
operation that actually cuts along a selected component can request the
full-coordinate mode and pay for its richer output.
